%----------------------------------------------------------------------------------------------------
% START OF CHAPTER
%----------------------------------------------------------------------------------------------------
\part{Numbers and Everyday Language} \label{M4-Numbers}

\chapter{Cardinal Numbers}

\section{0–20}
\begin{itemize}[nosep]
    \item \iti{zero, uno, due, tre, quattro, cinque, sei, sette, otto, nove, dieci}
    \item \iti{undici, dodici, tredici, quattordici, quindici, sedici}
    \item \iti{diciassette, diciotto, diciannove, venti}
\end{itemize}

\section{Tens}
\iti{venti, trenta, quaranta, cinquanta, sessanta, settanta, ottanta, novanta, cento}

\section{Key Spelling Rules}
\begin{enumerate}
    \item Drop final vowel of ten before \iti{uno} and \iti{otto}: \\
          \iti{venti + uno = ventuno, venti + otto = ventotto}
    \item \iti{Tre} gets accent at end of compound: \iti{ventitré, trentatré, quarantatré*}
    \item All compound numbers are \textbf{one word}: \\
          \iti{trecentosessantacinque} | Not \iti{\sout{trecento sessanta cinque}}
    \item \iti{Mille} → {mila} in plural: \iti{duemila, tremila}
    \item \iti{Cento} never changes: \iti{cento, duecento, trecento}
    \item \iti{Centouno} — hundreds keep their vowel before \iti{uno}
\end{enumerate}

\section{There is / There are}
\begin{itemize}[nosep]
    \item Singular: \itb{c'è} : \iti{c'è un problema}
    \item Plural: \itb{ci sono} : \iti{ci sono due problemi}
\end{itemize}

\chapter{Days, Months, Seasons and Time}

\section{Days of the Week}
Not capitalized: \iti{lunedì, martedì, mercoledì, giovedì, venerdì, sabato, domenica}

\begin{itemize}[nosep]
    \item Monday–Friday end in \itb{-dì} (accented)
    \item On a specific day: \iti{lunedì vado} (no article)
    \item Every week: \iti{il lunedì vado} (article = habitual)
    \item \iti{Domenica} is feminine; all others masculine
\end{itemize}

\section{Months}
Not capitalized, no article.
\begin{itemize}[nosep]
    \item \iti{gennaio, febbraio, marzo, aprile, maggio, giugno}
    \item \iti{luglio, agosto, settembre, ottobre, novembre, dicembre}
\end{itemize}

\section{Dates}
\itb{il + number + month} — no \iti{di} in between!

\begin{itemize}[nosep]
    \item \iti{il quindici marzo} | Not \iti{\sout{il quindici di marzo}}
    \item First of the month: \iti{il primo gennaio} (not \iti{il uno})
\end{itemize}

\section{Seasons}
Always \textit{in} without article: \iti{in primavera, in estate, in autunno, in inverno}

\section{Telling the Time}
\begin{table}[ht]
    \centering
    \begin{tabular}{r l l}
        .           & Rule                    & Example \\
           \hline
        Most hours  & \iti{sono le} + number & \iti{sono le tre} \\
        1 o'clock   & \iti{è l'una}          & \iti{è l'una}     \\
        Noon        & \iti{è mezzogiorno}    & -                 \\
        Midnight    & \iti{è mezzanotte}     & -                 \\
    \end{tabular}
\end{table}

\begin{itemize}[nosep]
    \item Past: \iti{sono le tre e venti} (twenty past three)
    \item Half: \iti{sono le tre e mezza} (half past three)
    \item Quarter: \itb{always} \iti{un quarto}: \iti{sono le tre e un quarto}
    \item To: \iti{sono le quattro meno un quarto} (quarter to four)
    \item At a time: \itb{alle} (\iti{alle otto}) | \itb{all'una} (at one)
\end{itemize}

\chapter{Common Expressions and Greetings}

\section{Greetings}
\begin{table}[ht]
    \centering
    \begin{tabular}{r l}
        Informal      & Formal \\
        \hline
        \iti{Ciao!} (hi/bye) & \iti{Buongiorno!} (good morning) \\
        \iti{Come stai?}     & \iti{Buonasera!} (good evening)  \\
        \iti{Come va?}       & \iti{Come sta?}                  \\
        \iti{Tutto bene?}    & \iti{Buonanotte!} (good night)   \\
    \end{tabular}
\end{table}

\section{Responses to "how are you?"}
\iti{Sto bene, grazie! | Non c'è male. | Così così. | Abbastanza bene. | Malissimo!}

\section{Introductions}
\begin{itemize}[nosep]
    \item \iti{Mi chiamo... | Come ti chiami? / Come si chiama? | Piacere! | Piacere mio!}
    \item \iti{Di dove sei? / Di dove è? | Sono di... | Vengo da...}
\end{itemize}

\section{Politeness}
\iti{Per favore / Per piacere | Grazie! | Prego! | Scusi! / Scusa! | Mi dispiace. | Non fa niente. | Permesso!}

\section{Key Expressions}
\begin{table}[ht]
    \centering
    \begin{tabular}{r l}
        Expression       & Meaning \\
        \hline
        \iti{Meno male!} & Thank goodness!   \\
        \iti{Meno male che...} & Thank goodness that... (always needs \iti{che} before a clause) \\
        \iti{Magari!}      & If only! / Maybe! \\
        \iti{Che peccato!} & What a shame!     \\
        \iti{Che bello!}   & How lovely!       \\
        \iti{Davvero?}     & Really?           \\
        \iti{Certo (che)!} & Of course!        \\
        \iti{Forse}        & Maybe/Perhaps     \\
        \iti{Allora...}    & So... / Well then... \\
        \iti{Dai!}         & Come on!          \\
        \iti{Basta!}       & Enough!           \\
        \iti{Andiamo!}     & Let's go!         \\
        \iti{Forza!}       & Come on! / Go for it! \\
        \iti{Già}          & Already / Indeed  \\
        \iti{Ancora}       & Still / Again     \\
    \end{tabular}
\end{table}

\section{Asking for Help}
\begin{itemize}[nosep]
    \item \iti{Non capisco. | Può ripetere? | Parla più lentamente, per favore.}
    \item \iti{Come si dice... in italiano? | Cosa significa...? | Non lo so.}
\end{itemize}

\section{At a Bar/Restaurant}
\iti{Vorrei... (I would like - polite) | Prendo... (I'll have) | Il conto, per favore. | Quant'è?}

%----------------------------------------------------------------------------------------------------
% END OF CHAPTER
%----------------------------------------------------------------------------------------------------